\documentclass[11pt]{article}
\usepackage[a4paper,margin=24mm]{geometry}
\usepackage{amsmath,amssymb,amsthm,mathtools,bm,mathrsfs}
\usepackage{booktabs,longtable,array}
\usepackage{enumitem}
\usepackage{xcolor}
\usepackage{microtype}
\usepackage[hidelinks]{hyperref}

\setlength{\emergencystretch}{3em}
\newtheorem{theorem}{Theorem}[section]
\newtheorem{proposition}[theorem]{Proposition}
\newtheorem{lemma}[theorem]{Lemma}
\newtheorem{corollary}[theorem]{Corollary}
\theoremstyle{definition}
\newtheorem{definition}[theorem]{Definition}
\newtheorem{gate}[theorem]{Fail-closed gate}
\theoremstyle{remark}
\newtheorem{remark}[theorem]{Remark}

\newcommand{\ii}{\mathrm i}
\newcommand{\dd}{\mathrm d}
\newcommand{\CC}{\mathbb C}
\newcommand{\RR}{\mathbb R}
\newcommand{\ZZ}{\mathbb Z}
\newcommand{\CP}{\mathbb {CP}}
\newcommand{\cO}{\mathcal O}
\newcommand{\cK}{\mathcal K}
\newcommand{\cL}{\mathcal L}
\newcommand{\cM}{\mathcal M}
\newcommand{\cC}{\mathcal C}
\newcommand{\cG}{\mathcal G}
\newcommand{\cE}{\mathcal E}
\newcommand{\cH}{\mathcal H}
\newcommand{\Tr}{\operatorname{Tr}}
\newcommand{\Ind}{\operatorname{Ind}}
\newcommand{\Ker}{\operatorname{Ker}}
\newcommand{\status}[1]{\textcolor{blue!60!black}{\textsf{#1}}}
\newcolumntype{L}[1]{>{\raggedright\arraybackslash}p{#1}}

\title{AP-E4 Same-Soliton Audit:\\
Conditional \(N=2\) Moduli-Space SQM, Families-Callias Determinant Line,\\
CPT/Anti-Canonical Polarization, and Gauge-Basic WZW Transgression}
\author{Route F research note}
\date{16 July 2026}

\begin{document}
\maketitle

\begin{abstract}
This note asks the stronger AP-E4 question that the independent vortex model
cannot answer: can a physical tangent fermion, its Callias determinant line,
and the level-two WZW line all be derived on the \emph{same} charged
two-colour \(B=1\) soliton?  The answer is a rigorous conditional theorem and
a rigorous underdetermination result, not a promotion claim.

At two-derivative adiabatic order, a one-multiplet tangent \(N=2\) SQM is
characterized at low energy by a smooth normalizable bosonic zero-mode
family, a uniform gap, a Fredholm fermion kernel
\(\cK_F\simeq T^{1,0}\CP^1=\cO(2)\), the \(N=2\) Ward identities, and the
Levi--Civita connection.  Exact microscopic supercharges are a sufficient
mechanism for these identities, but not a logical necessity because
accidental/emergent worldline supersymmetry is possible.  Our conservative
same-model audit therefore accepts either microscopic supercharges or an
independent derivation of all emergent Ward identities.  The present charged
two-colour proxy specifies a bosonic potential and selected Hessian blocks,
but no supersymmetry transformations and no Yukawa mass family.  Therefore a
gapped trivial fermion family and a topologically winding one are both
compatible with the current bosonic evidence.  The same-model tangent
fermion and Callias index are underdetermined.

For a declared rank-one asymptotic Callias eigenbundle over
\(S^2_\infty\times\CP^1\) with
\(c_1(\cE_+)=p x+q y\), the families boundary push-forward gives, in the
stated chirality convention,
\[
 \operatorname{rank}\Ind\mathscr D=\epsilon p,
 \qquad
 c_1(\det\Ind\mathscr D)=\epsilon pq\,y.
\]
Thus \(\epsilon p=1,q=2\) is the minimal rank-one template for
\(\cO(+2)\).  A gauge-invariant discrete Berry computation gives
\(c_1=1.999999999999981\), while a finite \(3\times3\) template/control kernel-family
Hamiltonian has residual below \(10^{-15}\) and unit nonzero gap.  These are
regressions of the conditional topology, not substitutes for the absent
four-dimensional Callias operator.

In a fixed Dolbeault polarization, CPT must use Serre duality:
\[
 E_{-}=K_{\CP^1}\otimes E_+^\vee=\cO(-4),
 \qquad (h^0,h^1)=(0,3).
\]
The naive raw inverse \(\cO(-2)\) has only one negative-chirality zero mode.
The necessary anti-canonical regulator factor is not derived by the current
mother model.

Finally, the degree-five mixed-WZW character is correctly transgressed along
the spatial three-cycle,
\[
 \widehat\cL_B=\int_{\Sigma_3}\operatorname{ev}^*\widehat\Omega_5
 \in\widehat H^2(\widetilde\cC_B;\ZZ).
\]
After an equivariant differential refinement and globally trivial vertical
holonomy establish gauge-basic descent, its pullback to the same soliton
moduli sphere is \(\cO(+2)\) exactly when
\(\int_{\Sigma_3\times\CP^1}\operatorname{ev}^*c(\widehat\Omega_5)=2\).
The factorized conditional map gives \(nB\deg(s)=2\), but the present proxy
has not identified its scalar \(S^2\) with the soliton collective sphere and
has not proved equivariant descent.  The AP-E3/AP-E4 composition gate and the
degree-one Route-E portal therefore remain closed.
\end{abstract}

\tableofcontents

\section{Scope, logic, and strongest defensible conclusion}

There are three spheres that must not be identified by notation alone:
\begin{enumerate}[label=(S\arabic*)]
\item the \(SU(2)\simeq S^3\) Skyrme field whose spatial degree is the
  baryon number \(B\);
\item the charged-scalar vacuum direction \(s\in S^2\) entering the mixed
  WZW target \(S^3\times S^2\);
\item a putative normalizable collective-coordinate sphere
  \(\cM_B\simeq\CP^1\) of the \emph{same} \(B=1\) soliton.
\end{enumerate}
An abstract diffeomorphism \(S^2\simeq\CP^1\) does not identify (S2) and
(S3).  Such an identification requires an explicit family of solutions and
an evaluation map.  The current nonlinear proxy establishes a product-target
\(B=1\) saddle and selected positive radial blocks.  It does not establish a
normalizable \(\CP^1\) degeneracy of that saddle, a four-dimensional
supersymmetric completion, or any fermionic fluctuation operator.

We therefore distinguish three kinds of statement:
\begin{description}
\item[Definition.] A construction that exists once its displayed geometric
  inputs are given, such as differential-cohomology fiber integration.
\item[Conditional theorem.] An implication with every sufficient audit
  hypothesis displayed explicitly; necessity is claimed only for the stated
  low-energy one-multiplet data, not for a microscopic supersymmetry mechanism.
\item[Same-model derivation.] A conclusion obtained from the actual charged
  two-colour microscopic action and its \(B=1\) solution.
\end{description}
Only the first two are achieved here.

\begin{theorem}[Strongest current AP-E4 conclusion]
The present charged two-colour bosonic data neither imply nor exclude a
physical tangent fermion.  They admit a precise conditional completion in
which the same-soliton Callias kernel is \(\cO(2)\), the \(B=+1\) SQM has
three positive-chirality ground states, the fixed-polarization \(B=-1\)
sector uses \(\cO(-4)\), and the spatially transgressed WZW line has degree
\(+2\).  None of these conditional inputs is selected by the current
fermion/Yukawa or gauge-equivariant data.  Consequently
\[
 \texttt{ap\_e3\_ap\_e4\_composition\_gate\_closed=false},
 \qquad
 \texttt{degree\_one\_route\_e\_portal\_built=false}.
\]
\end{theorem}

\section{Same-soliton collective-coordinate expansion}

\subsection{Configuration space and horizontal zero modes}

Let \(\mathscr S[\Phi,\Psi]\) be a hypothetical four-dimensional mother
action whose bosonic sector contains the charged two-colour candidate.  A
family of static unit-soliton solutions is a map
\begin{equation}
 \Phi_1:\Sigma_3\times\cM_1\longrightarrow\mathscr F,
 \qquad m\longmapsto\Phi_1(\cdot;m),
 \qquad B[\Phi_1(m)]=1,
 \label{eq:solution-family}
\end{equation}
where \(\mathscr F\) is the field-value space.  Gauge equivalence is removed
by a horizontal condition.  If \(m^a\) is a local complex coordinate and
\(\epsilon_a\) is the compensating gauge parameter, define
\begin{equation}
 \delta_a^H\Phi_1
 :=\partial_a\Phi_1-\delta_{\epsilon_a}\Phi_1,
 \qquad
 \mathscr G_\Phi(\delta_a^H\Phi_1,\delta_\lambda\Phi_1)=0
 \quad\text{for every }\lambda.
 \label{eq:horizontal-mode}
\end{equation}
The physical moduli metric is the Hessian kinetic pairing
\begin{equation}
 G_{a\bar b}(m)
 =\int_{\Sigma_3}\dd^3x\,
 \mathscr G_{\Phi_1(m)}
 \bigl(\delta_a^H\Phi_1,\delta_{\bar b}^H\Phi_1\bigr).
 \label{eq:l2-metric}
\end{equation}
Three requirements are separate: \(\delta_a^H\Phi_1\) must solve the
linearized equations, be \(L^2\)-normalizable, and have positive finite norm.
The current radial Hessian does not prove any of them for an orientational
mode.

Assume for the moment that \(\cM_1=\CP^1\).  On charts
\(U_N=\{w\}\), \(U_S=\{v\}\) with \(v=w^{-1}\), symmetry and direct chart
matching give
\begin{equation}
 G_{w\bar w}=\frac{C}{(1+|w|^2)^2},
 \qquad C>0,
 \qquad
 G_{v\bar v}|\dd v|^2=G_{w\bar w}|\dd w|^2.
 \label{eq:fs-metric}
\end{equation}
The number \(C\) must be computed from Eq.~\eqref{eq:l2-metric} for the same
charged soliton; borrowing it from a vortex does not close the gate.

\subsection{Fermion zero-mode family}

Let the quadratic fermionic action on \(\Phi_1(m)\) be
\begin{equation}
 S_F^{(2)}
 =\int\dd t\,\dd^3x\,
 \overline\Psi\bigl(\ii\partial_t-\mathscr D_m\bigr)\Psi,
 \qquad
 \mathscr D_m=\Gamma^i\nabla_i+\mathscr Y(\Phi_1(x;m)).
 \label{eq:fermion-family}
\end{equation}
Here the representation, gauge connection, grading, and Yukawa map
\(\mathscr Y\) are microscopic data.  If the kernel dimension is constant
and the rest of the spectrum has a uniform gap, the kernels form a smooth
Hermitian bundle
\begin{equation}
 \cK_F:=\bigsqcup_{m\in\cM_1}\Ker\mathscr D_m
 \longrightarrow\cM_1.
 \label{eq:kernel-bundle}
\end{equation}
For an orthonormal local frame \(\Xi_A(x;m)\), the adiabatic expansion
\(\Psi=\chi^A(t)\Xi_A(x;m(t))+\Psi_\perp\) gives the Berry connection
\begin{equation}
 (\mathcal A_a)^A{}_B
 =\int_{\Sigma_3}\dd^3x\,
 \Xi_A^\dagger\partial_a^H\Xi_B,
 \qquad
 D_t\chi^A=\dot\chi^A+(\mathcal A_a)^A{}_B\dot m^a\chi^B.
 \label{eq:berry-connection}
\end{equation}

\begin{theorem}[Low-energy characterization and sufficient microscopic closure]
Fix the two-derivative adiabatic truncation, demand one complex bosonic
coordinate, one complex Grassmann partner, no extra zero modes, and no
superpotential.  The low-energy tangent SQM requires the geometric/kernel
content (T1), (T3)--(T5) and the \(N=2\) Ward identities in (T2).  A sufficient
same-model microscopic closure is furnished when all of the following hold:
\begin{enumerate}[label=(T\arabic*)]
\item Eq.~\eqref{eq:solution-family} gives a smooth normalizable
  \(\CP^1\) family with positive metric and a uniform noncollective bosonic
  gap;
\item either the microscopic theory has two real supercharges preserved by
  the soliton and maps each \(\delta_a^H\Phi_1\) to a normalizable fermion
  zero mode, or an independent calculation derives the complete emergent
  worldline \(N=2\) Ward identities;
\item Eq.~\eqref{eq:fermion-family} is a uniformly Fredholm family and
  \(\cK_F\simeq T^{1,0}\cM_1\) as a Hermitian holomorphic line bundle;
\item the induced Berry connection under this isomorphism is the
  Levi--Civita connection, and the Ward identities fix the curvature
  four-fermion term;
\item all modes orthogonal to the collective multiplet remain uniformly
  gapped.
\end{enumerate}
\end{theorem}

\begin{proof}
At low energy, a physical collective coordinate requires (T1); a worldline
superpartner and its closure require the Ward-identity content of (T2);
tangent-valued transition functions require (T3); covariance of the
supercharge requires (T4); and a local adiabatic truncation requires (T5).
This does not make exact microscopic supercharges necessary: the same Ward
identities may arise accidentally after integrating out gapped modes.

For the stated sufficient closure, choose complex coordinate \(m\) and transport the normalized
kernel frame through the isomorphism \(\cK_F\simeq T^{1,0}\cM_1\).  Conditions
(T1)--(T5) give the standard covariant action
\begin{align}
 L_{N=2}
 &=G_{m\bar m}\dot m\dot{\bar m}
 +\ii G_{m\bar m}\bar\chi^{\bar m}D_t\chi^m
 -\frac12R_{m\bar m m\bar m}
 \chi^m\bar\chi^{\bar m}\chi^m\bar\chi^{\bar m},
 \label{eq:n2-action}\\
 D_t\chi^m
 &=\dot\chi^m+\Gamma^m{}_{mm}\dot m\chi^m.
 \label{eq:covariant-fermion}
\end{align}
Up to conventional factors, the transformations
\begin{equation}
 \delta m=\varepsilon\chi^m,
 \qquad
 \delta\chi^m=-\ii\bar\varepsilon\dot m
 -\Gamma^m{}_{mm}(\delta m)\chi^m
 \label{eq:susy-transform}
\end{equation}
and their conjugates close on time translation modulo the equations of
motion.  This is precisely the one-complex-dimensional tangent \(N=2\) sigma
model~\cite{Witten1982SameSoliton,AlvarezGaume1983SameSoliton}.
\end{proof}

On \(v=w^{-1}\), a tangent fermion must obey
\begin{equation}
 \chi^v=\frac{\partial v}{\partial w}\chi^w=-w^{-2}\chi^w.
 \label{eq:tangent-transition}
\end{equation}
The transition has degree two in the line-bundle convention, hence
\(\cK_F\simeq T^{1,0}\CP^1=\cO(2)\).  The verifier checks the metric,
fermion norm, and affine-connection transformation in both charts.

\begin{theorem}[Bosonic underdetermination]
The bosonic soliton family, its classical Hessian, and the mixed-WZW
coefficient do not determine \(\cK_F\) or even \(\operatorname{rank}\cK_F\).
\end{theorem}
\begin{proof}
Hold the bosonic action fixed.  Add a spectator vectorlike fermion with a
constant invertible gauge-invariant mass \(M\).  Its spatial operator has no
protected zero mode and its kernel bundle is zero.  Alternatively choose an
allowed Yukawa mass whose normalized asymptotic eigenspace winds over
\(S^2_\infty\times\cM_1\); the Callias boundary class below may then have
nonzero index and determinant Chern class.  These fermionic choices leave the
tree-level bosonic action and Hessian unchanged.  Therefore no function of
the bosonic data alone selects between them.
\end{proof}

The present charged two-colour theory is not supplied as a four-dimensional
supersymmetric action.  Emergent worldline supersymmetry is logically
possible, but would require the same equalities of metrics, Berry
connections, and interactions in (T1)--(T5); it cannot be inferred from an
abstract moduli-space complex structure.

\section{Families-Callias index and determinant line}

\subsection{Uniform Fredholm condition}

A Callias operator has the schematic graded form
\begin{equation}
 \mathscr D_m^+
 =D_{A(m)}^++\mathscr Y_m:
 H^1(\Sigma_3,E^+)\longrightarrow L^2(\Sigma_3,E^-).
 \label{eq:callias-operator}
\end{equation}
For every \(m\in\cM_1\), require constants \(R,\mu>0\), independent of
\(m\), such that outside the radius \(R\)
\begin{equation}
 \mathscr Y_m^\dagger\mathscr Y_m\ge\mu^2,
 \qquad
 \bigl\|[D_{A(m)},\mathscr Y_m]\bigr\|\le\mu^2-\delta
 \quad(\delta>0).
 \label{eq:uniform-callias}
\end{equation}
These hypotheses make the family Fredholm and define
\begin{equation}
 \Ind\mathscr D^+
 =[\Ker\mathscr D^+]-[\operatorname{coker}\mathscr D^+]
 \in K^0(\cM_1),
 \label{eq:index-bundle}
\end{equation}
with determinant line
\begin{equation}
 \lambda_{\rm C}
 :=\det\Ind\mathscr D^+
 =\det\Ker\mathscr D^+\otimes
 (\det\operatorname{coker}\mathscr D^+)^\vee.
 \label{eq:determinant-line}
\end{equation}
The Callias theorem and its families generalization reduce this class to the
positive-mass eigenbundle at spatial infinity
\cite{Callias1978SameSoliton,Kottke2011SameSoliton,VanDenDungen2024SameSoliton}.

\subsection{Boundary push-forward and the required class}

Let \(x\) and \(y\) be positive generators of
\(H^2(S^2_\infty;\ZZ)\) and \(H^2(\CP^1;\ZZ)\), respectively, normalized by
unit integral.  Consider the minimal rank-one product template
\begin{equation}
 c_1(\cE_+)=p x+q y,
 \qquad p,q\in\ZZ.
 \label{eq:positive-eigenbundle}
\end{equation}
Because \(x^2=y^2=0\),
\begin{equation}
 \operatorname{ch}(\cE_+)
 =1+p x+q y+p q\,x y.
 \label{eq:chern-character}
\end{equation}
Choose the Callias chirality/orientation sign \(\epsilon=\pm1\) so that the
boundary map is written
\begin{equation}
 \operatorname{ch}(\Ind\mathscr D^+)
 =\epsilon\int_{S^2_\infty}\operatorname{ch}(\cE_+).
 \label{eq:boundary-pushforward}
\end{equation}
This convention packages the possible overall sign from the spatial
orientation and grading.  Equations~\eqref{eq:chern-character} and
\eqref{eq:boundary-pushforward} give
\begin{equation}
 \operatorname{rank}\Ind\mathscr D^+=\epsilon p,
 \qquad
 c_1(\lambda_{\rm C})=\epsilon p q\,y.
 \label{eq:callias-rank-c1}
\end{equation}

\begin{proposition}[Minimal rank-one Callias target]
Within Eq.~\eqref{eq:positive-eigenbundle}, a single complex kernel with
\(\lambda_{\rm C}\simeq\cO(+2)\) exists if and only if
\begin{equation}
 \epsilon p=1,
 \qquad q=2,
 \label{eq:callias-target}
\end{equation}
assuming the cokernel vanishes and the kernel dimension stays constant.
\end{proposition}
\begin{proof}
Rank one requires \(\epsilon p=1\).  Substitution into
Eq.~\eqref{eq:callias-rank-c1} gives \(c_1=q y\), which equals \(2y\) if and
only if \(q=2\).  Constant kernel dimension and zero cokernel identify the
virtual determinant with the actual kernel line.
\end{proof}

The condition \(p=1,q=0\) has the same pointwise index but a trivial
determinant line.  A doubled pair \(q=+2\) and \(q=-2\) has vanishing total
determinant Chern class.  These are decisive negative controls: counting a
single zero mode at one moduli point does not compute a family Chern number.

\subsection{Berry formula and finite template/control regression}

When the kernel is an honest bundle with orthonormal frame \(\Xi_A\), its
determinant connection and Chern number are
\begin{align}
 \mathcal A&=\Tr\int_{\Sigma_3}\Xi^\dagger\dd_{\cM}\Xi,
 \qquad \mathcal F=\dd\mathcal A+\mathcal A\wedge\mathcal A,
 \label{eq:kernel-berry}\\
 c_1(\lambda_{\rm C})
 &=\frac{\ii}{2\pi}\int_{\CP^1}\Tr\mathcal F.
 \label{eq:berry-chern}
\end{align}
For numerical regression define the spin-\(k/2\) coherent vector
\begin{equation}
 u_k(\theta,\phi)_r
 =\binom{k}{r}^{1/2}
 \cos^{k-r}\frac\theta2\,
 \sin^r\frac\theta2\,e^{-\ii r\phi},
 \qquad 0\le r\le k.
 \label{eq:coherent-vector}
\end{equation}
It spans a smooth anti-Veronese Berry line of degree \(+k\).  With the
standard degree-\(k\) holomorphic structure this line is \(\cO(k)\); the
embedding into the trivial finite Hilbert bundle itself is not claimed to be
holomorphic.  On an oriented triangulation, the gauge-invariant Bargmann flux is
\begin{equation}
 \Phi_{abc}
 =\arg\!\left[
 \langle u_a,u_b\rangle
 \langle u_b,u_c\rangle
 \langle u_c,u_a\rangle\right],
 \qquad
 c_1^{\rm disc}=-\frac1{2\pi}\sum_{\triangle abc}\Phi_{abc}.
 \label{eq:discrete-chern}
\end{equation}
The minus sign is fixed by the anti-Hermitian convention
\(\mathcal A=u^\dagger\dd u\) and
\(c_1=(\ii/2\pi)\int\mathcal F\); it is not chosen after seeing the result.
The deterministic verifier uses 2208 triangles and obtains
\(c_1^{\rm disc}=1.999999999999981\) for \(k=2\).  It also tests
\(k=0,1,4\).

The finite matrix template/control family
\begin{equation}
 H_k(\theta,\phi)=\bm1_{k+1}-|u_k\rangle\langle u_k|
 \label{eq:finite-family}
\end{equation}
has an exact one-dimensional kernel, projector residual below machine
precision, and nonzero spectrum equal to one.  It verifies that a uniformly
gapped family with \(c_1=2\) is algebraically consistent.  It is not a
discretization of Eq.~\eqref{eq:callias-operator}; the actual charged-soliton
Yukawa data are still absent.

\begin{gate}[Callias status]
The formula \eqref{eq:callias-rank-c1}, its target class, and the Berry
regression are complete.  The same-model claims remain false because no
operator \(\mathscr D_m\), asymptotic eigenbundle \(\cE_+\), uniform mass gap,
or regulator coupling of \(\lambda_{\rm C}\) to the SQM wavefunction has been
derived.  In particular, a tangent kernel line does not automatically become
an additional coefficient line in the quantum Dolbeault complex.
\end{gate}

\section{CPT, Serre duality, and the anti-canonical sector}

\subsection{Why the naive inverse line is insufficient}

Quantize the conditional \(B=+1\) SQM in a fixed canonical Dolbeault
polarization with coefficient line
\begin{equation}
 E_+=\cO(2),
 \qquad
 Q_+=\sqrt2\,\bar\partial_{E_+},
 \qquad
 \cH_{0,+}=H^{0,*}(\CP^1,E_+).
 \label{eq:plus-polarization}
\end{equation}
Riemann--Roch gives
\begin{equation}
 (h^0(E_+),h^1(E_+))=(3,0).
 \label{eq:plus-count}
\end{equation}
A raw sign reversal of a degree-two Berry/WZW line would suggest
\(E_{\rm raw,-}=E_+^\vee=\cO(-2)\), but
\begin{equation}
 (h^0(\cO(-2)),h^1(\cO(-2)))=(0,1).
 \label{eq:naive-minus}
\end{equation}
Thus a signed integer \(+2\to-2\) alone does not construct the CPT triplet.

\subsection{Fixed-polarization anti-canonical map}

Serre duality on a complex curve states
\begin{equation}
 H^0(\cM,E)^*\simeq H^1(\cM,K_{\cM}\otimes E^\vee).
 \label{eq:serre-duality}
\end{equation}
For \(\cM=\CP^1\), \(K_{\cM}=\cO(-2)\).  Therefore the coefficient line that
realizes the conjugate triplet in the \emph{same} Dolbeault polarization is
\begin{equation}
 E_-=K_{\CP^1}\otimes E_+^\vee
 =\cO(-2)\otimes\cO(-2)=\cO(-4),
 \label{eq:anti-canonical-line}
\end{equation}
and
\begin{equation}
 (h^0(E_-),h^1(E_-))=(0,3),
 \qquad \operatorname{ind}\bar\partial_{E_-}=-3.
 \label{eq:minus-count}
\end{equation}

\begin{theorem}[CPT in a fixed Dolbeault polarization]
Let the \(B=+1\) zero modes be \(H^0(\CP^1,E_+)\).  An antiunitary CPT map
that reverses chirality while retaining the same displayed complex
polarization is implemented on zero modes by the Serre pairing if and only if
the \(B=-1\) coefficient line is
\(E_-=K_{\CP^1}\otimes E_+^\vee\).  For \(E_+=\cO(2)\), this is
\(\cO(-4)\), not \(\cO(-2)\).
\end{theorem}
\begin{proof}
The perfect Serre pairing
\begin{equation}
 H^0(E_+)\times H^1(K\otimes E_+^\vee)
 \longrightarrow H^1(K)\simeq\CC
 \label{eq:serre-pairing}
\end{equation}
gives an antilinear identification after choosing the Hermitian structure.
Conversely, a perfect fixed-polarization pairing with the positive sector for
all states identifies the negative cohomology with the Serre dual, fixing the
coefficient line up to holomorphic isomorphism.  Since
\(\operatorname{Pic}(\CP^1)=\ZZ\), its degree fixes it uniquely.
\end{proof}

An alternative description uses an anti-holomorphic involution and a
conjugate polarization.  When both sectors are rewritten in one fixed
polarization, the canonical factor in Eq.~\eqref{eq:anti-canonical-line}
reappears.  A physical charged two-colour CPT derivation must specify which
polarization and fermion-measure regulator are used.  The raw Callias or WZW
line only supplies the \(\cO(-2)\) part; the additional canonical factor is
independent UV data.  The mathematics of Eq.~\eqref{eq:anti-canonical-line}
is closed, but the same-model physical map is not.

\subsection{Finite SQM spectrum regression}

For the Fubini--Study metric
\(\dd s^2=C\,\dd w\dd\bar w/(1+|w|^2)^2\), the \(E_+=\cO(2)\) Dirac tower is
\begin{equation}
 \lambda_{n,\pm}
 =\pm\frac{2}{\sqrt C}\sqrt{n(n+3)},
 \qquad
 \operatorname{mult}(\lambda_{n,\pm})=2n+3,
 \qquad n\ge1.
 \label{eq:massive-spectrum}
\end{equation}
Serre duality gives the identical massive tower for the
\(E_-=\cO(-4)\) block, with the three zero modes in opposite chirality
\cite{BykovSmilga2023SameSoliton}.  The verifier builds four exact massive
levels, checks
\begin{equation}
 Q^2=(Q^\dagger)^2=0,
 \qquad \{\Gamma,D\}=0,
 \qquad H=\frac12\{Q,Q^\dagger\},
 \label{eq:sqm-algebra}
\end{equation}
and confirms indices \(+3\) and \(-3\).  This is a spectrum regression of
the conditional SQM, not evidence for its microscopic realization.

\section{Degree-five differential character and spatial transgression}

\subsection{The framed configuration-space line}

Let the mixed-WZW target be
\begin{equation}
 X=S^3\times S^2,
 \qquad
 \widehat\Omega_n
 =n\,\widehat u_3\smile\widehat u_2
 \in\widehat H^5(X;\ZZ),
 \label{eq:wzw-character}
\end{equation}
with normalized characteristic forms and classes
\begin{equation}
 \int_{S^3}\omega_3=1,
 \qquad
 \int_{S^2}\omega_2=1,
 \qquad
 c(\widehat\Omega_n)=n u_3u_2.
 \label{eq:normalizations}
\end{equation}
Differential characters and their fiber integration define holonomy without
choosing a five-dimensional extension
\cite{CheegerSimons1985SameSoliton,BaerBecker2014SameSoliton,
LeeOhmoriTachikawa2021SameSoliton}.

Let \(\widetilde\cC_B\) be the framed configuration component before the
gauge quotient and
\begin{equation}
 \operatorname{ev}:\Sigma_3\times\widetilde\cC_B\longrightarrow X,
 \qquad (x,\Phi)\longmapsto\Phi(x),
 \label{eq:evaluation-map}
\end{equation}
be evaluation.  Fiber integration gives the differential line class
\begin{equation}
 \widehat\cL_B
 :=\int_{\Sigma_3}\operatorname{ev}^*\widehat\Omega_n
 \in\widehat H^{5-3}(\widetilde\cC_B;\ZZ)
 =\widehat H^2(\widetilde\cC_B;\ZZ).
 \label{eq:spatial-transgression}
\end{equation}
Its curvature and characteristic class commute with fiber integration:
\begin{align}
 \frac{F_{\widehat\cL_B}}{2\pi\ii}
 &=\int_{\Sigma_3}\operatorname{ev}^*
   (n\omega_3\wedge\omega_2),
 \label{eq:transgressed-curvature}\\
 c_1(\widehat\cL_B)
 &=\int_{\Sigma_3}\operatorname{ev}^*(n u_3u_2).
 \label{eq:transgressed-c1}
\end{align}
This establishes the mathematically correct degree reduction.  It does not
yet give a line on the gauge quotient.

\subsection{Gauge-basic descent: local and global conditions}

Put \(\cC_B=\widetilde\cC_B/\cG\) and let
\(p:\widetilde\cC_B\to\cC_B\).  The exact descent statement is
\begin{equation}
 \widehat\cL_B\in\operatorname{im}\left[
 p^*:\widehat H^2(\cC_B;\ZZ)
 \longrightarrow\widehat H^2(\widetilde\cC_B;\ZZ)
 \right].
 \label{eq:descent-image}
\end{equation}
For a free smooth action, this is equivalent to a \(\cG\)-equivariant line
with connection whose curvature is basic and whose vertical action is
trivial.  Infinitesimally,
\begin{equation}
 \mathcal L_{\xi^\#}F=0,
 \qquad
 \iota_{\xi^\#}F=0,
 \qquad \xi\in\operatorname{Lie}\cG.
 \label{eq:local-basic}
\end{equation}
These local equations are not sufficient for disconnected/large gauge
transformations.  One must also require trivial holonomy on every vertical
loop and a trivial character on each stabilizer:
\begin{equation}
 \operatorname{Hol}_{\widehat\cL_B}(\gamma_g)=1,
 \qquad
 \rho_{\Phi}:\operatorname{Stab}_{\cG}(\Phi)\to U(1)
 \text{ is trivial}.
 \label{eq:global-basic}
\end{equation}

\begin{proposition}[Equivariant refinement criterion]
A sufficient intrinsic construction of Eq.~\eqref{eq:descent-image} is an
equivariant differential refinement
\begin{equation}
 \widehat\Omega_n^{\cG}\in\widehat H^5_{\cG}(X;\ZZ)
 \label{eq:equivariant-refinement}
\end{equation}
whose restriction is \(\widehat\Omega_n\), followed by equivariant fiber
integration.  It produces a line on the quotient stack.  Descent further to
the coarse quotient requires the stabilizer condition in
Eq.~\eqref{eq:global-basic}.
\end{proposition}

In physical terms, failure of Eq.~\eqref{eq:descent-image} is the residual
gauge anomaly of the WZW line.  Cancellation of local triangle anomalies in
a UV fermion table is necessary but does not alone compute the differential
holonomy for every large gauge loop.  The present model has not supplied
Eq.~\eqref{eq:equivariant-refinement} or Eq.~\eqref{eq:global-basic}; the
descent flag remains false.

\subsection{Pullback to the same soliton sphere}

Suppose the same charged soliton actually has a collective embedding
\begin{equation}
 \iota:\cM_B\simeq\CP^1\longrightarrow\cC_B.
 \label{eq:moduli-embedding}
\end{equation}
After descent, the line has degree
\begin{equation}
 k_{\rm eff}
 :=\int_{\CP^1}\iota^*c_1(\widehat\cL_B^{\rm quot})
 =\int_{\Sigma_3\times\CP^1}
  \operatorname{ev}_\iota^*(n u_3u_2).
 \label{eq:exact-pullback-integer}
\end{equation}

\begin{theorem}[Necessary and sufficient WZW pullback gate]
Assume gauge-basic descent and a fixed holomorphic structure on
\(\cM_B=\CP^1\).  The spatially transgressed line pulls back as \(\cO(+2)\)
if and only if the integer in Eq.~\eqref{eq:exact-pullback-integer} is \(+2\).
\end{theorem}
\begin{proof}
Naturality of characteristic classes and fiber integration gives
Eq.~\eqref{eq:exact-pullback-integer}.  Complex line bundles on \(\CP^1\)
are classified topologically and holomorphically by their degree because
\(H^2(\CP^1;\ZZ)=\ZZ\) and
\(\operatorname{Pic}(\CP^1)=\ZZ\)~\cite{Huybrechts2005SameSoliton}.
Hence degree \(+2\) is equivalent to \(\cO(+2)\).
\end{proof}

For the factorized conditional family
\begin{equation}
 \operatorname{ev}_\iota(x,m)=(g_B(x),s(m)),
 \qquad
 \int_{\Sigma_3}g_B^*u_3=B,
 \qquad
 \int_{\CP^1}s^*u_2=d,
 \label{eq:factorized-map}
\end{equation}
Eq.~\eqref{eq:exact-pullback-integer} becomes
\begin{equation}
 k_{\rm eff}=n B d.
 \label{eq:nBd}
\end{equation}
Thus \((n,B,d)=(2,+1,+1)\) gives \(+2\), whereas \(d=0\) gives zero and
\(n=1\) gives only one.  This is a useful sufficient ansatz and numerical
regression.  It is not a proof that the current scalar direction \(s\) is the
same normalizable collective coordinate as \(m\).

For \(B=-1\), the raw transgressed degree is \(-2\).  Section 4 shows that a
fixed-polarization CPT triplet instead needs total coefficient degree \(-4\).
The missing canonical \(-2\) must come from the CPT transformation of the
vacuum line/fermion measure, not from changing Eq.~\eqref{eq:nBd}.

\section{Composition theorem and fail-closed gates}

\begin{theorem}[Sufficient same-soliton AP-E3/AP-E4 composition criterion]
Within the field content and fixed-polarization construction specified in
this note, the AP-E3 level-two sector and AP-E4 SQM compose on one charged
two-colour \(B=1\) soliton when all of the following conservative audit
conditions are proved:
\begin{enumerate}[label=(C\arabic*)]
\item a stable \(B=1\) solution has a normalizable
  \(\cM_1=\CP^1\) family, finite positive \(L^2\) metric, and uniform
  noncollective bosonic gap;
\item the same four-dimensional action supplies either microscopic
  supercharges or independently derived emergent \(N=2\) Ward identities,
  together with a uniformly Fredholm Yukawa--Callias family;
\item its kernel bundle is \(T^{1,0}\CP^1\), with the required Ward identity
  and no extra zero modes;
\item the actual families index has rank one and determinant Chern number
  \(+2\), and the UV measure specifies whether that determinant line twists
  the wavefunctions;
\item the \(B=-1\) regulator derives
  \(E_-=K_{\CP^1}\otimes E_+^\vee=\cO(-4)\);
\item the WZW character admits an equivariant differential refinement and
  satisfies every local and global gauge-basic condition;
\item the same-soliton integral
  \(\int_{\Sigma_3\times\CP^1}\operatorname{ev}^*(2u_3u_2)=+2\);
\item no gauge projection, anomaly constraint, or lower-energy mixing removes
  the conditional three-state sector.
\end{enumerate}
\end{theorem}
\begin{proof}
Conditions (C1)--(C3) supply the tangent-SQM closure.  Condition (C4)
is the Callias determinant criterion and prevents a pointwise-index
substitution.  Condition (C5) is the CPT theorem.  Conditions (C6)--(C7) are
the necessary and sufficient differential-cohomology descent and pullback
criteria.  Condition (C8) is required for the Hilbert-space sector to survive.
Together they construct the claimed SQM.  These are sufficient closure
conditions for the present construction.  They are not asserted to exhaust
all possible emergent microscopic mechanisms.  With the current evidence,
each missing condition nevertheless invalidates the corresponding proposed
derivation step, so this composition has not been established.
\end{proof}

\begin{longtable}{L{0.43\textwidth}L{0.16\textwidth}L{0.31\textwidth}}
\toprule
Claim & Status & Reason\\
\midrule
\endhead
Low-energy tangent-SQM characterization & \status{closed} & required
low-energy data and a sufficient same-model audit are explicit\\
Bosonic underdetermination theorem & \status{closed} & constant-mass and
winding-Yukawa completions share the bosonic sector\\
Same charged-soliton \(\CP^1\) family & \status{open} & no solution family or
orientational \(L^2\) norm\\
Same-model \(N=2\) tangent fermion & \status{open} & no microscopic
supercharges, emergent Ward-identity derivation, or Yukawa operator\\
Conditional Callias \(c_1=+2\) template & \status{closed} & boundary
push-forward and Berry regression agree\\
Actual same-model Callias index/determinant & \status{open} & asymptotic mass
eigenbundle and uniform Fredholm gap absent\\
Mathematical anti-canonical CPT map & \status{closed} & Serre duality fixes
\(\cO(-4)\) in a fixed polarization\\
Physical same-model CPT regulator & \status{open} & canonical vacuum-line
factor not derived\\
Degree-five to degree-two fiber integration & \status{closed} & differential
cohomology defines Eq.~\eqref{eq:spatial-transgression}\\
Gauge-basic descent & \status{open} & no equivariant refinement or large-gauge
holonomy audit\\
Same-soliton pullback \(\cO(+2)\) & \status{open} & the exact integral is not
computed on a physical moduli family\\
AP-E3/AP-E4 composition & \status{open} & conditions (C1)--(C8) are not all
closed\\
Degree-one Route-E portal & \status{not built} & deliberately deferred until
composition closes\\
Physics promotion & \status{forbidden} & conditional mathematics is not a
microscopic derivation\\
\bottomrule
\end{longtable}

\section{Deterministic verification and reproducibility}

Run
\begin{verbatim}
PYTHONPYCACHEPREFIX=/tmp/python-cache \
python3 route_f/code/verify_ap_e4_same_soliton_callias_descent.py
\end{verbatim}
The program emits
\begin{verbatim}
route_f/output/ap_e4_same_soliton_callias_descent.json
route_f/output/ap_e4_same_soliton_callias_descent.md
\end{verbatim}
and hashes the TeX, bibliography, and verifier.  Its deterministic tests are:
\begin{enumerate}
\item CP1 metric, tangent norm, and Levi--Civita chart covariance;
\item gauge-invariant Berry-Chern numbers \(k=0,1,2,4\);
\item a uniformly gapped finite template/control kernel-family and its \(c_1=2\);
\item Callias boundary push-forward, \(q=0\), and doubled-family controls;
\item finite \(N=2\) matrices, nilpotency, chirality pairing, and the
  \(+3/-3\) Serre-dual indices;
\item the \(\cO(-2)\) one-state CPT negative control;
\item degree reduction \(5-3=2\) and the \(nBd\) transgression integer;
\item explicit fail-closed conditions for gauge descent and the same-model
  tangent theorem.
\end{enumerate}
The checks encode consequences of the analytic derivations.  They do not
prove the existence of the missing microscopic structures.

\section{Next calculations, in strict order}

Several routes remain viable, but they should be pursued without weakening
the gate:
\begin{enumerate}
\item \textbf{Supersymmetric same-model completion.}  Specify a
  four-dimensional charged two-colour action with two preserved real
  supercharges, solve its \(B=1\) family, and derive the coupled boson--fermion
  quadratic operator.  This is the direct route to (T1)--(T5).
\item \textbf{Non-supersymmetric Callias route.}  Even without exact
  supersymmetry, specify the physical quark/Yukawa operator, calculate its
  asymptotic positive-mass eigenbundle on
  \(S^2_\infty\times\cM_1\), and test whether
  \((\operatorname{rank},c_1)=(1,2)\).  A positive result gives a tangent
  Fermi bundle but still needs an emergent-SQM Ward-identity audit.
\item \textbf{Equivariant WZW route.}  Construct
  \(\widehat\Omega_2^{\cG}\), compute vertical holonomies for representatives
  of small and large gauge transformations, then evaluate
  Eq.~\eqref{eq:exact-pullback-integer} on the explicit same-soliton family.
\item \textbf{CPT regulator route.}  Compute the Bismut--Freed/Dai--Freed
  determinant-line transformation under \(B\mapsto-B\), and show whether the
  canonical factor \(K_{\CP^1}\) is supplied.  A raw \(-2\) result is not
  enough.
\end{enumerate}
Only after one coherent mother model closes all composition conditions should
a degree-one Route-E portal be written.  Until then the correct endpoint is
an explicit theorem of conditional composability plus a list of missing UV
data, which this note now provides.

\section{Machine-readable final flags}

The verifier records, among others,
\begin{align*}
 &\texttt{conditional\_composition\_sufficient\_criterion\_established=true},\\
 &\texttt{low\_energy\_tangent\_sqm\_characterization\_derived=true},\\
 &\texttt{sufficient\_same\_model\_audit\_conditions\_derived=true},\\
 &\texttt{underdetermination\_theorem\_applies=true},\\
 &\texttt{physical\_tangent\_fermion\_same\_model\_derived=false},\\
 &\texttt{callias\_determinant\_c1\_same\_model\_computed=false},\\
 &\texttt{conditional\_callias\_template\_c1\_plus\_2\_verified=true},\\
 &\texttt{cpt\_anti\_canonical\_map\_mathematically\_derived=true},\\
 &\texttt{cpt\_anti\_canonical\_map\_same\_model\_derived=false},\\
 &\texttt{wzw\_spatial\_fiber\_integration\_defined=true},\\
 &\texttt{wzw\_gauge\_basic\_descent\_same\_model\_proven=false},\\
 &\texttt{wzw\_pullback\_o2\_same\_model\_proven=false},\\
 &\texttt{ap\_e3\_ap\_e4\_composition\_gate\_closed=false},\\
 &\texttt{degree\_one\_route\_e\_portal\_built=false},\\
 &\texttt{physics\_promotion\_allowed=false}.
\end{align*}

\bibliographystyle{unsrt}
\bibliography{route_f/tex/ap_e4_same_soliton_callias_descent}

\end{document}
